\originalsection{相量、平均值与偏振的约定}
复相量并不是复数可测量场. 可测场为 $\ee(t)=\Re(\vct E_0e^{-i\omega t})$. 直接展开两个实部, 非零频率的周期平均满足
$\langle\Re(\vct a e^{-i\omega t})\times\Re(\vct b e^{-i\omega t})\rangle=\Re(\vct a\times\vct b^*)/2$.
因此 $\langle\vct S\rangle=\Re(\vct E_0\times\vct H_0^*)/2$, $\langle E^2\rangle=\vct E_0\cdot\vct E_0^*/2$. 共轭号不可省; 省去它会把不同相位误当同相.

写 $E_x=A\cos\omega t$, $E_y=B\cos(\omega t-\delta)$. 消去时间, $A,B\ne0$ 时有
\[
 (E_x/A)^2+(E_y/B)^2-2(E_x/A)(E_y/B)\cos\delta=\sin^2\delta.
\]
$\delta=0,\pi$ 退化为直线; $A=B$ 且 $\delta=\pm\pi/2$ 为圆; 一般为椭圆. 对圆偏振, 实场的恒定模长为 $A$, 相量却有两个幅值为 $A$ 的正交分量, 所以 $\vct E_0\cdot\vct E_0^*=2A^2$. 这正解释第17章圆偏振的平均能流为什么是 $A^2/Z_0$, 不再额外乘二分之一.

\originalsection{从纵场重建波导全场}
下面以真空、理想导体壁为例, 假定 $k_\perp^2=\omega^2/c^2-k_z^2\ne0$. 把麦克斯韦方程中的横向分量写出, 使用 $\partial_z=ik_z$, 例如
$\partial_yE_z-ik_zE_y=i\omega\mu_0H_x$、$ik_zE_x-\partial_xE_z=i\omega\mu_0H_y$.
与对应的安培方程联立, 消去另一横场后得
\begin{align*}
 \ee_\perp&=\frac{ik_z}{k_\perp^2}\nabla_\perp E_z
 -\frac{i\omega\mu_0}{k_\perp^2}\hat{\vct z}\times\nabla_\perp H_z,\\
 \vct H_\perp&=\frac{ik_z}{k_\perp^2}\nabla_\perp H_z
 +\frac{i\omega\eps}{k_\perp^2}\hat{\vct z}\times\nabla_\perp E_z.
\end{align*}
例如纯 TE 时 $E_z=0$, 横电场全部由 $H_z$ 的横向梯度给出. 壁上的切向方向是 $\hat{\vct z}\times\nn$, 投影上式得到 $E_\mathrm{wall,t}\propto\partial_nH_z$; 故 TE 用 Neumann 条件. TM 时 $H_z=0$, 纵电场本身是壁切向场, 必须令 $E_z=0$; 沿壁的切向微分也随之为零, 保证横电场切向分量为零. 因而纵场的标量本征问题和全部矢量边界是一致的.

TEM 时 $E_z=H_z=0$ 且 $k_\perp=0$, 不能在上式中直接代零, 因其推导除过 $k_\perp^2$. 另从横向方程得 $\nabla_\perp\times\ee_\perp=0$ 和 $\nabla_\perp\cdot\ee_\perp=0$, 故 $\ee_\perp=-\nabla_\perp V_\perp$, $\nabla_\perp^2V_\perp=0$. 单连通空心波导只有一条导体边界, 所有边界同一常势; 唯一性使 $V_\perp$ 恒定, 电磁场皆零. 同轴结构有两条导体边界, 两常势可不同, 因而存在非零TEM场. 这证明了有无TEM的几何原因, 不只是背诵一个模式名称.
